\section{文档范围与源码基线}
\lead{NDNSF、NDNSF-UAV、NDNSF-DI、NDNSF-Repo 完整设计视图}
本文覆盖四个模块的职责、主要接口、控制与数据流程、安全、状态、恢复、配置、部署和实现边界。按子系统解释设计，不逐行转录源代码。代码标识、协议字段和文件名保留原名，其余叙述使用中文。

\begin{description}
\item[版本] R3；当前设计记录已核对源码行为与明确限制，补充具体接口、状态及例子。目标设计保留独立基线，另列五项 PLANNED 改进。
\item[基线] \SnapshotDate，\code{Experimental} 分支。采样时提交为 \texttt{\SourceBaseline}；逐文件摘要、工作树差异与本地源码快照见 \code{source-baseline.json}。
\item[事实口径] 主要流程核对源码；已有接口与已接通流程分别说明。历史测试仅作为检索入口，不冒充本轮结果。
\item[版本维护] 本目录完整纳入 Git。两份正文独立保存；修改目标设计不自动修改当前设计或代码。
\end{description}

\begin{tabularx}{\linewidth}{@{}p{30mm}X@{}}
\toprule
内容组 & 覆盖内容 \\
\midrule
整体与 Core & 所有者关系、调用、协作、数据传输、并发与状态。\\
授权与安全 & 请求级保密、版本、授予/撤回、状态传播与更新影响。\\
Repo & 对象/制品、持久化、传输、目录、副本与恢复。\\
DI & 应用接口、规划、安全、原生执行、生成/KV、部署与扩展。\\
UAV & 容器、飞控、任务、操作权限、视频、录像及多视角识别。\\
交付与证据 & 跨模块流程、运行环境、实现状态、源码索引及维护。\\
\bottomrule
\end{tabularx}

\callout{关键现状}{当前已验证的原生请求入口是
Runtime::open → User::prepare → PreparedModel → NativeInferenceClient，
并由 C++ 过程选择器覆盖 ACK、Selection、Provider、Response、会话和清理。
没有 Runtime preparation/request contract 的兼容构造仍明确返回
\code{NATIVE_REQUEST_PIPELINE_NOT_READY}；这条限制不能被已验证的 PreparedModel 路径掩盖。}

\newpage

\section{系统组成与职责边界}
\lead{Core 提供通用机制；应用模块拥有领域语义}
\begin{center}
\begin{tikzpicture}[box/.style={draw=navy,rounded corners,fill=lightblue,align=center,text width=43mm,minimum height=15mm,font=\small},>=Stealth]
\node[box] (u) at (0,0) {NDNSF-UAV\\飞控、任务、媒体与识别};
\node[box] (d) at (5.2,0) {NDNSF-DI\\模型、规划、执行与生成};
\node[box] (r) at (10.4,0) {NDNSF-Repo\\命名对象与制品存储};
\node[box,text width=147mm] (c) at (5.2,-2) {NDNSF Core\\调用、授权、协作、状态、命名数据与流};
\node[box,text width=147mm] (n) at (5.2,-3.8) {NDN 与运行依赖\\ndn-cxx / NFD / ndn-svs / NAC-ABE / 密码与存储库};
\draw[->,thick] (u)--(c);
\draw[->,thick] (d)--(c);
\draw[->,thick] (r)--(c);
\draw[->,thick] (c)--(n);
\end{tikzpicture}
\end{center}
UAV 可调用 Repo 保存录像，也可调用分析服务；DI 可调用 Repo 发布、复用或获取模型制品。
这些是应用间依赖，不意味着 UAV 飞行状态、模型图或 Repo 目录语义进入 Core。[C1, R1, D1, U1]

\begin{tabularx}{\linewidth}{@{}p{28mm}X@{}}
\toprule
模块 & 代码位置与所有权 \\
\midrule
Core & \path{ndn-service-framework/}；通用协议与 C++ runtime。\\
UAV & \path{NDNSF-UAV-APP/}；drone、ground-station、shared 及工具/配置。\\
DI & \path{NDNSF-DistributedInference/}；C++ 角色执行与迁移中的原生请求端、Python 应用和策略层。\\
Repo & \path{NDNSF-DistributedRepo/}；C++ 库/服务与 Python 制品编排。\\
Python 绑定 & \path{pythonWrapper/}；Core 绑定与应用辅助层，不等价于全部逻辑均在 C++。\\
\bottomrule
\end{tabularx}

\callout{边界规则}{Core 可以传递、认证和保存应用定义的 payload，但不解释 waypoint、tensor、model shard、codec 或 object class 的领域含义。}

\newpage

\section{Core：运行角色、命名与启动}
\lead{身份、服务与进程容器相互关联，但不是同一个对象}
\begin{description}
\item[ServiceController] 维护授权策略和撤权状态，提供签名权限响应及 PolicyStatus，持有 NAC-ABE Authority。
\item[ServiceUser] 发布请求、收集 ACK、选择参与者、发送 Selection、处理 Response；维护请求、token、版本与超时状态。
\item[ServiceProvider] 注册服务和 ACK/执行处理器；响应候选发现、验证选择、运行任务、发布结果。
\item[ServiceContainer] 同一进程组合 Controller/User/Provider 角色和本地服务，管理启动、停止及生命周期 hook。
\item[LocalServiceRegistry] 提供可信进程内调用；不是网络调用模式，不增加线上的权限属性或服务选择语义。[C1, C2]
\end{description}

\subsection{名字与消息}
服务通过 NDN 名称寻址；Request、ACK、Selection、Response 由 Core 的类型和命名 helper
编码。Provider 身份、服务名、requestId 与消息类别共同约束匹配；具体文法以
\code{NDNSFMessages.*} 和 \code{utils.*} 为准，不由应用拼装新的旁路协议。

通用服务可由多个 Provider 提供；应用需要指向一个物理节点时可使用 Provider 专属服务名，
或在已经完成 token bootstrap 后使用 Targeted。名字可达、服务被发现和当前获准执行是不同条件。

\subsection{启动与安全准备}
加载配置、身份/证书、信任模式，建立 Face 和服务注册，再获取权限和 ABE 材料。
构造期异步获取 DKEY 不能阻塞整个应用运行；但 bootstrap-pending 不构成执行授权。
Controller 的启动含独立 Face 的公参可达性探测，避免仅以内进程命中宣称 NFD 路径可用。[C2, S3]

配置集中在 ConfigManager、PolicyParser 及应用配置文件；服务容器只负责组合，硬件连接、
模型加载和数据库路径由各应用模块决定。


\subsection{从配置到可调用服务}
以一个进程同时拥有 User 和 Provider 为例：容器身份不是服务名；先配置 Face、证书和 Controller 前缀，再把 /example/echo 注册给 Provider，并保存 scoped registration。容器 start 后事件循环才处理发现与执行。User 取得 requestId 只表示本地请求建立。完整签名、句柄保存与 stop 顺序分别见 AC-02 至 AC-04。

\newpage

\section{Core：普通调用与 Targeted 调用}
\lead{先完成发现与选择，再执行；已知 Provider 可使用受保护的短路径}
\begin{enumerate}
\item User 为新 invocation 分配 requestId、token 和时间预算，发布 Request。
\item 可服务且具有相应权限的 Provider 返回 ACK，携带 Provider token、能力和应用自定义报价。
\item ACK 窗口结束后，User 从候选中选择一个或多个 Provider。策略可为 first-responding、random、all-selected 或自定义。
\item User 发布绑定实际候选与 token 的 Selection；未选中节点不能据 Request 自行开始执行。
\item Provider 验证选择、权限与请求绑定，调用处理器；User 验证返回的签名、密钥和 requestId，再触发结果回调。[C1, C3]
\end{enumerate}

\begin{center}
\begin{tikzpicture}[>=Stealth,font=\small]
\node (u) at (0,0) {User};
\node (p) at (9,0) {Provider 候选};
\draw[gray] (0,-.3)--(0,-3.3); \draw[gray] (9,-.3)--(9,-3.3);
\draw[->,navy] (0,-.7)--node[above]{Request：发现描述}(9,-.7);
\draw[->,navy] (9,-1.4)--node[above]{ACK：候选、token 与能力}(0,-1.4);
\draw[->,navy] (0,-2.1)--node[above]{Selection：精确选择及接收者密钥封装}(9,-2.1);
\draw[->,navy] (9,-2.8)--node[above]{Response：执行结果或终态}(0,-2.8);
\end{tikzpicture}
\end{center}

\subsection{Targeted 的适用条件}
已知 Provider 且已通过正常认证流程取得一次性 token pair 后，后续调用可省略 ACK/Selection。
权限、Provider 绑定、请求/响应认证和 token 防重放仍然存在。token 缺失、过期或耗尽时，
应走相应 bootstrap/错误处理，不能退化成无认证的直发。[C3]

\subsection{失败的归属}
Nack、没有可用 ACK、负 ACK、选择失败、执行错误、响应验证失败与超时是不同边界。
发布成功不等于执行成功；普通 ACK 也不天然表示独占资源预留。


\subsection{应用怎样判断完成}
例如请求预算为 3000 ms，应用将 requestId 与界面操作关联，等待 ResponseHandler 或 TimeoutHandler。ACK 阶段不能显示执行完成，也不应启动另一个重复业务请求。普通请求与 Targeted 的并列声明和最小调用见 AC-03；涉及 token 的短路径仍要通过同一服务权限检查。

\newpage

\section{Core：协作规划与选择事务}
\lead{不可变 ACK 快照与一次计划提交把应用规划接入同一 invocation}
\code{DEFERRED} 协作先发起通用请求，再收集 \code{ACK_CLOSED} 快照；应用根据快照决定
角色、参与者、依赖和密钥作用域。规划完成后调用同一次协作的 \code{commit_plan}。
\code{PREPLANNED} 保留请求前已知计划的兼容入口，两者不能被混成一次动态重开 ACK 窗口。[C3, D2]

\subsection{Core 负责的契约}
角色标识、精确参与者绑定、依赖名称/作用域、请求与 attempt 关联、选择消息分发及终态处理，
由通用协作契约承载。模型切分可行性、UAV 任务分区或 Repo 副本评分由上层判断。

\code{GenericSelectionTxnStore} 保存认证后的选择上下文与应用提交 blob；
\code{OpaqueSelectionParticipant} 提供 prepare、onCommitted、onAborted。
事务关联 Provider boot epoch、token、可选 lease、截止时间、payload digest 和存储 key epoch；
重放已提交记录也检查当前运行代次。[C4]

\subsection{带资源预留的选择}
R1 的 DI reservation 是独立协商能力，不能赋予所有普通 ACK 相同语义。
启用时 Provider 先完成有界资源预留，再给 positive ACK；选择后提交/准备，
未选中或迟到候选收到 \code{NOT_SELECTED}。资源最终由本地 lease 到期收回。

选中 source role 本地就绪即可运行；有依赖的角色还要获得其直接前驱的认证数据。
不把旧兼容 ReadySet/ExecutionActivate 接口解释为所有当前路径都需要的全局启动屏障。[D2]

\callout{完成边界}{计划提交、Provider 接受、资源准备完成、角色执行完成、终端结果被接受是不同状态；文档和应用状态显示应分别记录。}


\subsection{一个有依赖的协作计划}
设 A 在 P1 执行，B 在 P2 执行且消费 A 的输出。ACK\_CLOSED 固定本次候选，应用将 A/B 的参与者、A 到 B 的命名依赖及终端规则写入计划，再以相同 requestId 和闭包摘要提交。Core 接受提交后，P2 仍需验证选择并获得 A 的实际输出。提交失败不能靠提前执行 B 补救。AC-05 将 User 提交与 Provider 选择事务的两个 owner 分开说明。

\newpage

\section{Core：对象传输、连续流与流式调用}
\lead{按数据语义选择入口，不仅按字节数选择}
\begin{tabularx}{\linewidth}{@{}p{30mm}X@{}}
\toprule
数据形态 & 当前接口与用途 \\
\midrule
完整命名对象 & \code{publishLargeNamed}、\code{fetchLarge}、\code{fetchLargeExact}；确定名字和版本的模型、tensor bundle、manifest、录像文件。\\
连续发布流 & \code{StreamPublisher}、\code{PredictiveStreamSubscriber}、\code{StreamInfo/StreamChunk}；不断产生的媒体、遥测与状态。\\
单次调用的事件流 & \code{RequestServiceStreaming}、\code{StreamedResponseWriter}；一个 invocation 内的有界事件、取消和终止，如生成 token。\\
\bottomrule
\end{tabularx}

\subsection{完整对象}
请求中可携带小型引用；实际字节按精确 Data 名获取并验证。大响应的请求级保密路径按段
AEAD，必须具有相应的请求密钥/绑定。Repo 增加持久存储和副本语义，不建立另一套无命名大 payload 通道。[C5]

\subsection{连续流}
通用流管理会话身份、cursor、名称映射、序列缺口、重排、重复、保留窗口、预测请求与背压。
当前 Stream 还提供对不透明字节的有界 FEC 修复能力；媒体帧类型、编码参数和画面呈现仍归 UAV。
流恢复必须受期限/保留范围限制，不能无限阻塞后续交付。[C6]

\subsection{调用事件流}
事件绑定 service、request、Provider、ControllerVersion 和 generation 等上下文。
有界队列、事件保留和重试、transcript/terminal 校验、取消与 deadline 共同控制生命周期。
终态只被接受一次，晚到事件不得复活已经结束的请求。普通调用与 Targeted 均须满足各自授权前提。[C7]

\callout{选择例子}{直播使用连续流；已完成录像使用命名对象；一个推理请求逐 token 输出使用调用事件流。三者可以组合，但不能互相冒充。}


\subsection{对应到代码类型}
完整对象的结果是可验证的全部字节；连续流的 StreamPublisher::push 输入已经签名的 Data，返回 void；调用事件流的 writer.publish 输入 EventT，返回 bool，finish 则可能抛异常。三种 API 的参数和返回值不能混用，分别见 AC-06、AC-07。

\newpage

\section{Core：能力、并发、持久状态与观测}
\lead{公共状态表达事实，上层策略解释业务意义}
\begin{description}
\item[能力与发现] 通用 capability/runtime hint 表达可用性、负载、队列、资源和 drain 状态；应用字段以带 schema 的 payload 附着。可达不代表可以接收新任务。
\item[准入与租约] ProviderAdmissionLeaseTable 和 ExecutionLease 维护有期限的准入/执行状态；request、service、Provider、资源绑定及消费/释放检查防止跨请求复用。
\item[服务生命周期] 注册可使用作用域句柄，关闭后按注册代次阻断晚到回调。一个 Provider host 可共享物理资源 lease 状态，避免多个服务各自重复预订同一槽位。
\item[操作状态] 通用状态区分排队、执行、等待输入、完成、失败、取消和过期；结果以类型化 metadata 或 DataProductReference 表达。[C2, C8]
\end{description}

\subsection{线程与终态}
Face 的 I/O 线程、ACK/handler 工作池以及应用后台线程承担不同工作。
pending 请求和资源表由相应锁/代次保护；取消和关闭要清理绑定，不能仅删除 GUI 条目。
Core 的事件回调、请求 deadline 与应用长生命周期任务仍分开拥有。

\subsection{持久化不等于恢复完成}
ControllerGenerationStore 管理持久授权版本；RuntimeStatusStore 可选保存签名状态并重新验证；
选择事务有自己的 WAL/存储密钥。DI 的请求 journal、UAV 的 mission ledger、Repo 的数据库
是上层持久状态，不能用一个文件存在代替全部恢复握手。[C4, S3, S12]

\subsection{可观测性}
NetworkTelemetry、TimelineTrace 和应用 diagnostics 提供事件与时序信息。
应区分排队/网络/验证/装配/计算/呈现时间；成功标记、源码 hash 和原始运行目录需要关联。
单项测试或某个统计计数不自动构成端到端性能或安全资格结论。


\subsection{一次执行槽位的生命周期}
同一物理设备的冲突资源通过 conflictKeys 与 proof 绑定。租约从 Prepared，经 commit 成为 Committed，validateAndActivate 再核对 \code{requester/request/service/plan} 进入 Executing，结束后 release；未选预留可 abort，到期可 cleanup。Provider 重启后的 providerEpoch 变化会使旧代租约不可照用。AC-09 给出完整创建、提交、激活、释放、查询签名与 reasonCode，观测计数按这些状态解释。

\newpage

\section{Core：身份验证与请求级保密}
\lead{签名验证、服务授权和内容保密解决不同问题}
\subsection{身份与信任}
证书与 Trust Schema 确认 Data 的签名身份及名称是否符合信任规则。
Controller 的授权策略决定身份能否使用或提供服务；签名有效本身不等于拥有服务权限。
加密证书与签名证书可分开配置，但须满足同一身份及运行要求；NAC-ABE 使用的加密证书需要兼容其密钥封装。[C2, C5]

\subsection{默认 V2 请求级保密路径}
\begin{enumerate}
\item 初始 Request 的加密发现信息面向获授权的 Provider 类别，使用户不必预先知道具体执行者。
\item 选择后，为本次调用准备独立的输入/响应密钥；输入作为 User 签名的精确命名 Data，
在输入密钥下加密。
\item SelectionKeyEnvelope 使用选中 Provider 的 RSA-OAEP 封装传递
\code{K_input} 与 \code{K_response}，绑定接收者证书及请求上下文。
\item 响应使用 AES-GCM；AAD 包含 service、requestId、ControllerVersion、
attempt、证书摘要等绑定。分段或事件使用不同 nonce，nonce registry 拒绝复用。
\item 终态/失效路径释放或清零请求秘密；大响应仍需逐段验证及请求绑定，不能退回共享服务响应钥匙。[C5, S7]
\end{enumerate}

\subsection{ABE 与其他密钥的范围}
ABE 公参和 DKEY 用于属性策略下的授权/保密；请求 AES 密钥用于本次输入/输出；
HybridMessageCrypto 仍用于其他服务级、协作或 token 场景。
DI artifact grant、UAV 视频内容钥匙和操作权 lease 是各自应用契约，不与这些字段混成一个 epoch。

\subsection{证书和公参可达性}
安全 bootstrap 需要实际名称路由、Controller/Authority 可达和签名验证。
NFD 前缀注册、forwarding hint 或本地缓存只影响数据如何到达，
不能取代身份和内容验证；失败要按证书、公参、权限、状态或内容层定位。


\subsection{按攻击对象选择检查}
签名有效但没有服务属性的 User 仍应被拒绝；具有服务属性但未被选定的 Provider 不能获得本次输入密钥；得到模型文件的节点也不因此具有特定请求的受保护模型 grant。调用者依次核对身份、服务授权、request/selection 绑定与领域 grant，不把四个问题缩成一个 authenticated 标志。

\newpage

\section{Core：授权版本与权威}
\lead{版本推进、密钥轮换和权限变化是三个不同判断}
\begin{tabularx}{\linewidth}{@{}p{36mm}X@{}}
\toprule
对象 & 当前语义 \\
\midrule
ControllerVersion & 两个非零 uint64：generation timestamp 与 epoch；先比较 generation，再比较 epoch。Controller 的同一全局版本用于各服务的 PolicyStatus。[S1, S2]\\
PolicyStatusData & 按服务返回的签名状态：ControllerVersion、有效期、策略摘要、ABE 公参名字和摘要、Controller 证书、相关撤权记录。[S2]\\
ABE generation & 以实际公参 Data name + digest 区分。grant-only 保持不变；withdrawal 轮换。节点据此判断是否清空共享 NAC 缓存。[S3, S6]\\
PermissionResponse & 面向目标身份的权限表，携带 policy epoch。节点替换本地记录并比较权限内容，决定是否需要目标 DKEY 刷新。[S5]\\
DKEY & 当前 KP-ABE 实现按身份持有完整属性策略的解密密钥，并非每次新增一个服务就附加一个独立属性密钥。[S3, S4]\\
\bottomrule
\end{tabularx}

\subsection{不能仅从字段推导的结论}
\code{advanceAuthorizationEpoch()} 同时设置 \code{m_policyEpoch} 与
\code{m_requiredKeyEpoch} 为新 epoch；这不表示每次调用都进行了 ABE 全局轮换。
运行时另行比较公参身份。[S3]

\code{getPolicyStatus()} 对每个服务装入相同的 ControllerVersion，再过滤服务级撤权记录；身份级和证书级撤权也被加入状态。当前结构不是独立的“每身份权限修订号”。[S2]

该函数的 \code{policyDigest} 来源于配置文件 SHA-256；不能把它当作完整在线 grant/revoke 变更集摘要，也不能仅凭其不变就判断某节点不受影响。[S2]

\subsection{权威边界}
经认证的业务消息所携带的新版本只能触发状态拉取。安装状态必须验证 Controller 签名、服务、有效期及版本绑定；旧版本被拒绝，同版本但内容冲突也被拒绝。状态与刷新协调器先在副本上验证，再一起安装。[S6, S8, S9]

\callout{不能混用的两个结论}{“ABE 参数未变”意味着不必全局重发 ABE 密钥；它本身不证明该服务所有在途请求仍满足新的授权状态。}

\newpage

\section{Core：在线授予权限}
\lead{保留 ABE 参数，仅替换目标身份的完整策略}
\begin{enumerate}
\item 应用或控制入口调用 \code{ServiceController::grant(identity, service, attribute)}。User 使用 \code{/PERMISSION/...}，Provider 使用 \code{/SERVICE/...}。
\item 验证目标及 Controller 状态；若此前有未完成的撤权轮换，先恢复该轮换，失败则拒绝本次 grant。证书撤销不可由 grant 绕过。
\item 更新目标身份的服务授权和属性集合，移除匹配的服务撤权；显式重新授权可解除该身份的身份级撤权。
\item 替换目标身份的完整有效属性策略，重建查找表，再推进全局授权 epoch。
\item Authority 在该身份下一次请求时生成替换 DKEY。普通 grant-only 不调用全局 ABE 参数轮换，也不遍历其他身份重发策略。[S4]
\end{enumerate}

\subsection{目标身份如何获得新增权限}
权限发现需要应用层续期或显式 \code{fetchPermissionsFromController()}。仅知道新 PolicyStatus，并不自动等价于获取新增权限；User/Provider 对权限响应和状态响应分别处理。

\code{applyPermissionResponse()} 比较更新前后的权限内容：User 比较服务集合；Provider 比较 Provider/service 记录集合。仅内容确实改变时，标记相关服务等待 DKEY 更新。随后安装状态时，若 ABE 参数未变且存在待刷新标记，则执行 DKEY-only refresh。[S5, S6]

同一 ControllerVersion 的多个服务响应合并为一次身份 DKEY 更新；状态先到、权限后到的反向顺序也由同版本安装路径处理。相同代次刷新期间，旧密钥保留到替换密钥安装。[S6, S10]

\subsection{其他身份会发生什么}
没有权限内容变化且 ABE 公参身份不变时，无关身份不会仅因为更高 epoch 而请求替换 DKEY。不过，Provider 接受已安装状态的版本推进后会重新拉取自身权限表；运行状态失效也仍然可能发生。详见“Core：授权变化的运行影响”一节。[S5--S7]

\callout{本流程的限定范围}{上述不轮换结论限定于普通 grant-only。若 grant 入口先恢复了此前失败的撤权轮换，此次调用包含那次轮换的影响，不能归类为纯目标身份更新。}


\subsection{单项输入与策略物化}
给身份 U 新增服务 S 时，调用 grant(U, S, attribute)，控制器向已有服务/属性集合插入这一项，再从有效集合重建 ABE 策略。应用无需提交 U 的完整权限表。相同授权已存在可以返回 false；证书已撤销等拒绝也可以返回 false，所以返回值不能单独区分所有原因。已有待轮换状态时先恢复该状态，grant-only 的局部更新描述不能覆盖这个例外。见 AC-08。

\newpage

\section{Core：在线撤回权限}
\lead{撤回权限，同时切换后续加密使用的 ABE 代次}
\begin{enumerate}
\item \code{revoke(target)} 支持身份、证书或特定身份的服务权限撤回；先验证目标并处理已有 pending rotation。
\item 将撤权目标加入集合，持久推进 ControllerVersion。持久化失败时恢复本次内存修改。
\item 轮换 Authority 的 ABE 主密钥/公共参数代次；公参版本由 generation 与 epoch 派生。
\item 对已知身份遍历剩余有效属性：先移除旧策略；无剩余权限则不再设置策略，否则设置完整过滤后策略。客户端后续获取对应 DKEY。
\item 后续 PolicyStatus 返回当前版本、当前 ABE 公参身份和相关撤权记录；User/Provider 拉取、验证并安装后，触发授权与缓存处理。[S3, S8]
\end{enumerate}

\subsection{失败处理顺序}
授权 epoch 已持久化而 ABE 轮换抛错时，不回退到旧版本。撤权保留，设置
\code{m_abeRotationPending}，调用返回失败。下一次 revoke（包括同目标重试）或 grant 先执行恢复；恢复会预留更高 epoch，再轮换，以免改写已经发布的同版本状态。[S3]

这是 Controller 本地失败处理，不表示离线或尚未更新的远端节点已经获知撤权。Controller 成功、远端发现状态、客户端重新取钥、受保护操作拒绝旧授权，是不同观察点。

\subsection{密码学与运行时边界}
新代次下加密的数据不能用旧代次 DKEY 正常解密。已经交付的旧明文、旧密钥或旧代次密文不因此被追溯收回。运行中的操作还依赖本地版本门、请求绑定与失效处理，不能仅用“已换钥”描述全部撤权效果。

在节点安装不同公参身份后，会清理本地共享 NAC Consumer 缓存、获取新 DKEY，并清理/刷新 Producer 的公参状态。这个共享 ABE 范围大于单个服务的本地 token 范围。[S6, S7, S10]

\callout{证据边界}{本文件描述代码顺序和拒绝路径；没有给出撤权全网生效时延、网络分区下的完成保证或本轮运行测试结论。}

\newpage

\section{Core：版本发现与传播}
\lead{当前通过 NDN 拉取状态；没有独立的 Controller SVS epoch 公告}
\begin{tabularx}{\linewidth}{@{}p{38mm}X@{}}
\toprule
触发方式 & 当前行为 \\
\midrule
权限获取/更新 & 权限路径触发相应 PolicyStatus 获取；新增权限依赖应用续期或显式 refetch。[S5]\\
经认证的新版本 hint & 协调器检查 hint 是否可信、有效且更新；合并正在进行的拉取，按目标版本确认签名状态。[S9]\\
临近过期刷新 & 默认按当前状态到期时间安排，携带当前精确版本；旧精确版本请求本身不能发现新 epoch。[S11]\\
周期重新检查 & \code{NDNSF_POLICY_REVALIDATION_PERIOD_MS} 大于 0 时使用固定周期，并发出不指定版本的 fresh 状态请求。默认值 0，关闭此周期模式。[S11]\\
可选状态恢复 & 恢复持久状态先重新验证，再向 Controller 做在线确认；这不是对离线最新状态的证明。[S12]\\
\bottomrule
\end{tabularx}

\subsection{Interest 与 Data 契约}
不指定版本的请求前缀为
\begin{quote}\small
\path{<controller>/NDNSF/POLICY-STATUS/<service>}
\end{quote}
客户端设置 \code{CanBePrefix=true} 和 \code{MustBeFresh=true}；
Controller 用当前精确版本名返回签名 Data，\code{FreshnessPeriod=0}。
指定版本时客户端关闭前缀匹配；Controller 若已不提供该版本则不回答。
客户端验证 Data 名字、内部版本和期望版本的一致性。[S8]

\subsection{SVS 在当前设计中的位置}
ServiceController 的当前实现未建立 epoch 专用 SVS publisher/subscriber。
框架其他业务使用 SVS，不代表授权更新已经具备相同公告路径。
本文件及首版目标设计均不把“SVS 通知后拉取状态”方案写成已实现能力。

\callout{时效边界}{仅推进 Controller epoch 不会主动唤醒所有节点。周期模式、显式拉取和消息 hint 的触发条件决定何时发现变化；本轮没有测量发现或完成时间，也不承诺所有节点立即同时切换。}

\newpage

\section{Core：授权变化的运行影响}
\lead{本地已实现部分区分，但尚非“无关节点完全无感”}
设 Provider A 获得服务 S 的新权限，Provider B 的权限内容没有变化，
ABE 参数仍为 P，Controller 版本由 V 变为 V+1。以下描述 B 或 User
\textbf{实际接受服务 S 的新状态之后} 的代码分支，而非自动广播后的全网结果。

\begin{tabularx}{\linewidth}{@{}p{40mm}X@{}}
\toprule
状态或动作 & 当前结果 \\
\midrule
B 的 ABE DKEY & 比较权限内容无变化，且公参身份未变时，不执行 DKEY-only refresh。Authority 也没有替换 B 的策略。[S4--S6]\\
B 的权限表 & Provider 从已有状态推进版本时，触发自身权限 refetch；LocalMock 跳过这一网络动作。[S7]\\
服务 S 的本地状态 & 新版本安装调用 \code{invalidateControllerScopedCaches}，并未先判断变更是否仅属于另一个身份。[S6]\\
User 的在途调用 & 同服务 stream 被拒绝；普通调用保留 pending/deadline，但清除请求秘密、绑定、token 等，让终态/超时路径继续负责结束。[S7]\\
Provider 的在途状态 & 清理匹配服务的 pending 请求、token、stream 绑定和相关发布状态；不是按 grant 的目标身份过滤。[S7]\\
其他服务 T & 本次服务 S 的失效按服务匹配。但各服务状态携带同一全局 ControllerVersion；若稍后接受 T 的更新状态，T 也会进入自己的失效路径。[S2, S6, S7]\\
\bottomrule
\end{tabularx}

\subsection{当前设计的准确结论}
可以说：普通 grant-only 只替换目标身份的 ABE 策略；无关身份不因 epoch
变化而重新取 DKEY。

不能说：普通 grant-only 只影响目标 Provider，其他节点无须拉取或更新任何状态、
所有在途调用不受影响。这超过当前运行状态失效逻辑所支持的结论。

\subsection{当前行为与已接受目标}
本页记录全局 ControllerVersion 的基线行为。目标 TG-01 已接受细化影响范围的方向，
字段线格式、通知承载及轮换粒度仍待后续 Spec 定稿；NDN-SVS 是可选通知方案，
尚未据此建立当前实现的广播保证。目标还须规定版本兼容、离线补齐和在途请求处置。

\newpage

\section{Repo：架构与对象模型}
\lead{命名对象存储与领域无关；C++ 库和 Python 编排各有边界}
\begin{description}
\item[RepoCore] 不依赖 NDNSF 网络调用的进程内存储门面；提供 put/get/list/remove、manifest、Data packet、容量和目录操作。
\item[RepoNode] 将 RepoCore 挂接到 ServiceProvider 或 LocalServiceRegistry，注册存储操作服务，管理请求响应及 operation 状态。
\item[RepoClient] 提供直接、进程内和网络 helper；Python 的 py\_repoclient 进一步提供副本、目录和制品会话编排。
\item[存储后端] 普通对象/packet 的持久后端与大制品的 PayloadStore/MetadataStore 接口分别承载不同物化方式。[R1--R4]
\end{description}

\begin{tabularx}{\linewidth}{@{}p{35mm}X@{}}
\toprule
数据对象 & 关键身份与含义 \\
\midrule
RepoObjectManifest & objectName、objectType、SHA-256、size、segmentCount、replicationFactor、replicaNodes、policyEpoch、以及由 Core 持有的 protected durable publication identity（publication、protection epoch、opaque key-reference、ciphertext manifest digest、serving locator）和生命周期 metadata；Repo 只保存并返回这些非秘密字段。\\
已签名 Data packet & 按原始 NDN 名称与 wire bytes 保存；保持发布者身份，不能因进入 Repo 就改成 Repo 内容签名。\\
ArtifactReference & 不可变制品身份：逻辑名、内容摘要、大小、root manifest 名称、发布者与策略 epoch；不是只有路径。\\
目录条目 & 逻辑对象、manifest 摘要、来源 Repo、状态、目录版本及副本位置；不等于全部段字节。\\
\bottomrule
\end{tabularx}

\subsection{控制面与数据面}
默认服务前缀为 \code{/NDNSF/DistributedRepo}。CAPABILITY、INSERT、MANIFEST、
INVENTORY、STATUS、CACHE\_STATUS、CATALOG 系列和 DELETE 等属于控制操作；
具体 C++/Python 可用操作以各自注册表为准，不能假定两个入口完全对等。
数据面仍通过原始精确名称获取分段 Data，forwarding hint 可帮助定位副本。[R2, R7]

\callout{通用性}{Repo 不解释模型层、视频帧和张量。payload 是否加密、如何解码和何时可以用于业务，由发布者、消费者及应用权限机制决定。}


\subsection{对象引用的阅读顺序}
先看 objectName 定位逻辑对象，再用 manifest 中的摘要、大小、segmentCount 和有序 packet 名验证内容，最后核对 replica metadata 的实际可读性。模型文件和录像 packet 可以共用 Repo 存储，但前者的模型语义由 DI 解释，后者的签名原始 wire 由 UAV 验证。AC-10 给出两份 packet 的保存/获取例子。

\newpage

\section{Repo：持久化、写入与原子可见性}
\lead{先提交权威存储，再更新缓存；收到字节不等于发布完成}
\subsection{普通对象与 packet 存储}
部署中的普通 Repo-node 对象路径使用 SQLite 权威存储，叠加按字节预算限制的 LRU 热点缓存。
写入先提交数据库，再纳入缓存；读取 miss 回到数据库；删除/覆盖先完成权威事务。
超过缓存预算的条目旁路缓存；缓存和计数随进程消失，持久对象不会因缓存清空而被视为丢失。
in-memory backend 是测试替身，不应当成受支持的持久部署模式。[R3]

\subsection{大制品后端}
\code{PayloadStore} 管理分范围写入、verified ranges、flush/finalize/abort；
\code{MetadataStore} 管理生命周期事件。\code{RepositoryStoreFacade} 结合后端所有权 lease，
防止多个 owner 无约束地修改同一存储目录。\code{FilesystemArtifactStore}
和本地 Python artifact backend 支持文件/staging 物化，与普通 SQLite 对象 API 不能混称为一个实现。[R4]

\subsection{网络写入流程}
\begin{enumerate}
\item 发布者确定精确制品身份和 manifest，公开供选中 Repo 拉取的数据源。
\item 通过 NDNSF 协作发现容量/能力，固定副本选择和 operationId。
\item 各 Repo 获取、校验并持久化数据，产生绑定 operation、artifact 和 Repo 身份的提交 receipt。
\item 客户端验证收到的 receipt 与已提交的选择集合及内容身份一致，满足要求后才返回发布完成。
\end{enumerate}
分段收到、哈希验证、持久提交和对外可见分别记录。副本不足或 receipt 不匹配不能用“部分上传成功”包装成完整发布。[R5, R6]

\callout{范围区别}{普通存储 helper 的重试/分段适配器与 artifact-manifest-v2 的会话生命周期不是同一层；不能把某一条本地路径通过当作所有网络后端都通过。}


\subsection{成功信号属于哪个阶段}
普通 put 的本地返回、制品 transfer 完成、commit 产生已验证 receipt、目录观察到对象是不同事件。断点状态只能支持继续接收，不能让未验证片段直接对外宣称完整文件。文档中的原子可见性指对应存储入口的提交边界，不向所有 Python 编排、缓存和跨 Repo 目录扩展为一个分布式原子事务。

\newpage

\section{Repo：大制品传输与可恢复会话}
\lead{公开 API 隐藏 packet 编排，保留精确身份和提交检查}
\begin{tabularx}{\linewidth}{@{}p{46mm}X@{}}
\toprule
接口 & 语义 \\
\midrule
publish\_file / publish\_file\_async & 校验源文件身份，发布到要求的副本数，返回不可变引用及副本结果。\\
fetch\_file / fetch\_file\_async & 按引用获取并校验文件，完成后物化到目标位置。\\
begin\_upload → upload\_file → commit & 显式上传会话，区分传输与提交。\\
begin\_fetch → transfer → commit & 显式获取会话，区分下载与最终可见文件。\\
\code{abort(preserve_progress=True)} & 保留已验证 staging，允许在相同身份与恢复条件下续传。\\
\bottomrule
\end{tabularx}

\subsection{清单、分段与背压}
artifact-manifest-v2 将 root manifest、清单页和实际 payload 分开，避免把海量段列表塞进控制消息。
AdaptiveArtifactTransfer 跟踪 in-flight、重传、重复和 verification backlog；
收到段先离开 Interest 窗口，但在验证/持久化完成前继续施加背压。
恢复会话绑定 artifact、root digest、chunk/packet 参数和 lease；不能把相同文件名当作可续传证明。[R4--R6]

\subsection{当前网络入口的限制}
\code{CollaborationArtifactApiBackend} 承载网络制品发布；当前
\code{_NetworkPublishDriver.transfer()} 明确拒绝 TARGETED 整制品发布，
要求 Collaboration。API 中存在 Targeted 控制选项，不表示每个具体 backend 的所有操作都支持它。[R6]

默认路径按能力与资源边界选择副本，并验证选中集合对应的持久 receipt。
本地单副本 backend 用于本地运行/开发，不代表已经经历真实 NDN 多副本传输。

\callout{恢复原则}{只复用经过验证的精确进度；过期 lease、不同 digest、不同分段参数或错误 operation 绑定应由会话检查拒绝，不能静默接着写入另一个对象。}


\subsection{断点下载的操作身份}
同一 artifact 下载到不同目标路径可能形成不同 operation。fetch\_file 默认 resume=True、verify=True、replace=False；先 begin\_fetch，再 transfer，最后 commit。异常时依据 resume 决定是否保留进度。AC-11、AC-12 展开 descriptor/result、receipt、副本计数、kwargs 和底层 receive/markVerified 的区别，避免把进度条完成当成文件提交。

\newpage

\section{Repo：目录、副本、保留与恢复}
\lead{目录负责定位和生命周期；对象字节仍由数据面获取}
每个 Repo 有本地目录，Python 控制编排还提供目录交换、查询和修复工具。
系统没有把单一集中 NameNode 设为所有对象读取的必要前提；各入口实际是否启用 peer 同步，
取决于所选服务和配置，而不是仅因为存在 catalog 类型就自动发生。[R1, R2, R7]

\subsection{目录接口}
\begin{description}
\item[CATALOG\_STATUS] 返回 Repo 模式、目录 epoch、对象数量和副本接收能力。
\item[SNAPSHOT / DELTA] 提供全量恢复或指定 epoch 之后的变化；目录 sequence 不等于 Controller 授权 epoch。
\item[C++ LOOKUP] catalogLookup 仅接受精确 objectName 并返回一个条目，没有组合过滤参数。
\item[Python QUERY] metadata 过滤和目录编排属于 Python 路径；不能归入 C++ LOOKUP 的能力。
\end{description}

\subsection{副本选择与确认}
候选 ACK 中的 freeBytes、mode、负载等用于 placement；选中不代表已经保存。
策略需满足复制数量/容量约束，真正成功由验证后的写入 receipt 确认。
已知副本位置可优化请求，但仍保留名称、签名、manifest 和内容摘要检查。

\subsection{删除、TTL 与修复}
删除以 tombstone/目录状态表达，避免旧 AVAILABLE 记录无条件恢复对象。
对象可带 \code{ttlMs} 与 \code{repairAllowed}；过期或不允许修复的对象不应作为新副本修复来源。
恢复机制需区分 metadata 恢复、字节验证和实际 replica 完成，不能仅依据目录列表计副本数。[R7]

\subsection{应用集成}
UAV 的内嵌 Repo 为加密录像和原始 packet 留存提供存储；DI 使用公开 artifact API 传输模型制品。
独立 Repo-node、应用内 RepoCore 和 Python 网络制品后端是不同部署形态；文档中的能力应绑定具体入口。


\subsection{一次目录读不等于副本验收}
catalogLookup 在目录锁外读取 store manifest，再在锁内附加本 Repo 的 epoch，因此不是存储与目录的全局一致快照。missing/store 错误沿路径传播，不是统一返回空集合。P1 的 epoch 不能作为 P2 的 delta 游标；目录条目指向 P2 也仍需对实际字节取数验证。精确查询能力和 Python 编排边界见 AC-13。

\newpage

\section{DI：应用入口与当前实现层次}
\lead{原生 PreparedModel 是请求入口；Python 只做薄适配和编排}
\code{ndnsf_distributed_inference.api} 导出 Runtime、User、PreparedModel、句柄和
Provider 的 pybind 对象；\code{api/_async.py} 只转换 completion/reader 到 asyncio。
旧 \code{app_sdk}/\code{APPClient} 编排仍保留为兼容 caller，不能把兼容入口的状态
外推为 PreparedModel 已覆盖的全部能力。[D1]

\subsection{两类应用请求}
\begin{description}
\item[模型意图请求] 以 ModelRef、task、input 和 timeout 提交；AutomaticPlanningCoordinator
在 ACK 窗口关闭后产生图与分割候选、验证 Provider offer、决定计划并提交。
\item[显式预规划请求] 使用 RequestableDeployment 或既定 plan；保留兼容 API，
不应与无需应用先指定 Provider/role 的延迟规划语义混用。
\end{description}

\subsection{C++ 已有的能力}
Runtime/User/PreparedModel 将已验证模型包、输入投影和请求句柄绑定到
NativeInferenceClient；NativeExecutionPlan、NativeProviderSession/Runtime、
ProviderRoleWorker、模型 runner、依赖传输、生成协调器及 Provider host 负责后续阶段。
原生规划、grant、ONNX、tokenizer、会话和请求生命周期由 Spec185 的 C++ 证据逐项覆盖。[D3--D6]

\subsection{当前路径与限制}
PreparedModel 的 request/run 会把准备包、输入、选项、绝对 deadline 和策略交给
NativeInferenceClient，过程资格已由同一候选的 C++ selector 验证。
直接构造没有 preparation/request contract 的 NativeInferenceClient 仍保留
\code{NATIVE_REQUEST_PIPELINE_NOT_READY}，用于明确区分兼容/组件构造与生产 prepared 路径。
Python subinterpreter、wheel packaging 和 Spec184 外部模型/实验资格仍未观测。[D7]

\callout{目标与基线}{目标 TG-02 继续要求同一原生实现配可选 Python 外壳；当前
PreparedModel 路径已经是该方向的已实现部分，兼容入口和未观测外部资格仍单独记录。}


\subsection{如何选择入口并解释失败}
需要现在的应用编排时定位 app\_sdk/client.py 的 APPClient.request；需要检查原生迁移则定位 \code{NativeInferenceClient.cpp} 的 dispatchOperation。模型/任务/输入是业务值，native handle 是异步操作句柄，两者不能替代。AC-14 说明输入形式、默认预算、局部等待和取消；全局实现状态集中见第 31 章，目标所有权迁移见 TG-02。

\newpage

\section{DI：输入、模型图、候选与放置}
\lead{先冻结真实观测，再做模型相关的规划}
\begin{enumerate}
\item 应用输入经过 adapter 预处理，绑定模型内容/语义摘要与输入类型；大输入使用命名引用。
\item 以 DEFERRED 协作获取不可变 ACK\_CLOSED，绑定 request、attempt 和总 deadline。
\item 由固定 adapter 描述模型并检查图，枚举有界、图有效的 split candidates。
\item 从真实 ACK payload 构造 ProviderOfferV3；检查身份、offer 签名、Core 认证来源、模型/图绑定与新鲜度。
\item 放置策略评估候选的完整 role/rank、backend/device、内存、队列和驻留证据；
独立验证策略输出，不把策略返回值直接作为执行授权。[D2, D3]
\end{enumerate}

\subsection{驻留证据与准备}
loaded、assembled、canonical 与 cold 表示不同准备成本。
\code{has_model} 不能证明 exact residency；加载证据还需匹配进程/boot、runtime、
设备、artifact、graph、role 和有效期。具有所需服务能力但尚未加载的 Provider，
可以在契约允许时接受带准备的执行；不能通过伪造“已驻留”来获得资格。

\subsection{策略扩展}
ModelFamilyAdapter 拥有图/切分/物化语义；ModelPlacementStrategy 选择 Provider 与设备。
PreSplitFirst 系列优先复用合法、精确匹配的候选与驻留证据，再考虑准备成本。
策略受候选数量、期限和验证器限制，不能扩大自己的权限或绕过 Controller。

\subsection{当前原生迁移}
NativeRequestPreparation、NativeOfferAdmission、NativeV3Placement 已有对应组件和检查；
当前真实数据来源端口及 requester 主链仍需接线。原生组件的定向测试不能替代
模型意图经过网络 ACK 到终端结果的整体验收。[D3, D7]


\subsection{用两段模型读懂数据结构}
线性图分为 A、B 两段，A 输出必须匹配 B 输入，这一分解形成 candidate。认证报价分别给出 P1/P2 的角色、设备与资源，placement 才提出 A→P1、B→P2 的 proposal；validator 再校验它。模型描述、图快照、candidate、offer、proposal 依次收窄自由度，不是同一份 JSON 的五个别名。AC-15 给出精确复用与需准备路径的例子，AC-16 解释发布后如何重绑身份。

\newpage

\section{DI：计划封装、执行授权与制品}
\lead{计划完整性、权限与制品保密分别绑定}
计划包含模型/图身份、角色与 rank、Provider assignment、依赖边、设备和执行约束。
V3 的 proposal、plan core、Provider grant view 与 selection projection 分开表示不同
可见范围；规范编码和摘要用于验证绑定，不能以临时 JSON 字符串或目录顺序代替身份。[D2, D4]

\subsection{从决策到选择}
\begin{enumerate}
\item 验证完整候选与 assignment，确定需要的可信物化和制品。
\item 通过负责该模型的 adapter/assembler 生成合法切片或装配产物。
\item 使用公开 Repo API 发布/确保制品可获取，保留内容与 manifest 身份。
\item 封装不可变计划与授权绑定，向选中 Provider 提交其最小必要投影。
\item Provider 验证 selection、计划、角色、lease、证书/代次与制品绑定后再进入准备/执行。
\end{enumerate}

\subsection{原生安全组件}
NativeGrantClient/Verifier、NativePlanSealer、NativeProtectedArtifactStore、
NativeProtectedProvider 和 ProtectedRuntime 分担 grant、验证、受保护制品和明文生命周期。
ProtectedRuntime 的 lease/zeroization 组件用于显式受保护路径；不能因此推导所有普通
模型加载或默认 Python 路径都已启用同等的模型制品加密。[D4]

\subsection{三个不同的安全问题}
Core 的请求保密保护调用输入/输出；DI 的模型/计划身份保证执行的是被指定的内容；
受保护 artifact grant 控制特定制品的解密与运行时明文持有。三者必须分别说明，
不能用“使用了 ABE”概括所有模型权限。

\callout{当前限制}{Spec182 的封装、grant 和准备 owner 虽已有实现片段，真实请求端接线与全部验收尚未闭合；接口和类型存在不等于默认网络路径已使用它。}


\subsection{五种摘要不互换}
ACK 闭包摘要绑定观察到的候选；模型/图摘要绑定被检查内容；artifact root 绑定发布的业务制品；plan core 摘要绑定执行内容；grant/projection 绑定授权视图和具体参与者。发布切片后必须以真实 root/recipe 再绑定角色，再封存计划，不能改下载 URL 却保留旧驻留证明。AC-16 按 prepare、inspect、ensure、bind、seal、grant、project、encode 说明前后数据。

\newpage

\section{DI：原生角色执行与依赖数据流}
\lead{角色执行在本地就绪后推进，依赖边通过认证命名对象连接}
\begin{tabularx}{\linewidth}{@{}p{43mm}X@{}}
\toprule
组件 & 主要职责 \\
\midrule
NativeExecutionPlan & 角色、依赖边、张量与执行身份；NativeProviderAssignment 指定实际 owner。\\
NativeProviderSession & 将计划、assignment、DependencyIo 和 runner factory 组合成会话。\\
NativeProviderRuntime & 维护 runner 注册与 worker，安排角色执行。\\
ProviderRoleWorker & 等待/预取直接输入、运行角色、按依赖边发布输出；管理局部状态和缓存。\\
NativeProviderHandler & 将 Core CollaborationContext 接到原生 session，验证执行绑定、readiness 与授权。\\
\bottomrule
\end{tabularx}

\subsection{输入与输出}
TensorBundleCodec 定义张量 bundle 编解码；模型 adapter/runner 验证 dtype、shape、
名字、图输入/输出及必要的边界转换。DependencyIo 把数据流边映射为精确命名对象，
在 producer/consumer/request/attempt/plan 等绑定下获取和发布。
同机角色可使用本地 I/O；网络角色通过 NDNSF 的数据传输接口协作。[D5]

\subsection{调度与资源}
DependencyWaitScheduler/AsyncDataflowRuntime 负责在依赖可用时推进，避免把等待网络的任务
无限占住执行线程。CollectiveRuntime、ProviderGroupCoordinator 等组件表达 group/rank 协调；
这些类型和执行路径的存在不能自动作为多 GPU、扩展性或吞吐测量结果。

\subsection{模型运行器}
OnnxRuntimeModelRunner 以及 ONNX/Qwen/YOLO adapter 处理各自的模型格式和运行契约。
设备、CUDA/CPU execution provider、runner 身份与 runtime ABI 必须按实际配置核对；
不能把允许 GPU 的配置当成实际使用 GPU。输出发布和 terminal role 接受仍由当前执行绑定约束。[D5]


\subsection{A 到 B 的实际执行}
Session 组合 plan、assignment、DependencyIo、runnerFactory 与 worker 数。A 的初始输入可由 initialInputs 提供；worker 运行后通过已绑定的依赖 scope 发布 TensorBundle，B 在对应前驱输入就绪后运行。executeRoleAsync 返回 future，异常需由调用者取回并收尾；发布一个中间张量不等于终端结果被 User 接受。session、I/O 和 runner 生命周期须覆盖 future 的工作期，见 AC-17。

\newpage

\section{DI：逐 token 生成、KV 与会话延续}
\lead{不可变模型缓存与带身份的推理状态缓存分开管理}
\subsection{生成循环}
NativeEpochCoordinator 组织每轮生成：读取输入/反馈、校验 lineage、执行角色、
抽取 logits、按生成配置采样 token、发布 token/终止事件，并处理取消与截止时间。
QwenGenerationSession、GenerationEpochLineage 和 DecodeStateIdentity 记录模型、前缀、
position、epoch 及角色状态之间的关系。[D6]

\subsection{流与终态}
生成事件经 Core 的调用事件流传输；stable text decoding 需要处理 token/字节边界、
特殊 token、最终 flush 和不能稳定输出的前缀。EOS、stop sequence、max tokens、失败、
取消和 deadline 是不同结束原因。最终结果的归属不能由任意中间角色声明。

\subsection{KV 复用}
KV 不是普通不可变模型 artifact。可复用状态必须匹配模型内容与语义、adapter、
runner、split、tokenizer、精确 prefix/长度、position、层范围、precision/layout、
runtime ABI、Requester/安全域及有效期。过期、驱逐、重启或 pin 丢失不能继续作为命中证据。
Provider session 提供 host/GPU pause、prefetch、pin、release 等显式状态操作。[D5, D6]

\subsection{跨轮对话}
ConversationStateBinding 和 NativeConversationCoordinator 用于绑定会话/轮次及状态来源，
Python app journal 也保存请求和恢复记录。延续不是重放上一轮的未经验证对象；
旧 attempt、晚到结果和已取消状态必须经过相应门禁。

\callout{验收边界}{当前已有 tokenizer、epoch 和 conversation 组件；Spec182 的 sampling parity、
stable epoch emission、journal/continuation 与完整原生 stream 接线仍保留未完成验收项。
不把某些部件测试通过写成完整生成路径已迁移。}


\subsection{一次生成与下一轮对话}
实际入口 \code{runNativeEpochCoordinator} 逐 role/epoch 调用 runner.run，模型状态机执行 Selecting→Prefilling→Decoding→Draining→Completed。attempt 是替换尝试，streamEpoch 是流实例，context epoch 是对话轮次，三者不按每个 token 同步递增。AC-18 给出方法状态表、两 token 示例、KV pin/迁移/释放以及 beginTurn→prepareCheckpoint→commitTurn 的完整阅读主线。

\newpage

\section{DI：部署、扩展与管理接口}
\lead{服务注册、制品准备、执行与管理动作各有生命周期}
\begin{description}
\item[应用/部署层] app\_sdk.application、deployment、controller 和 runtime\_journal 管理定义、
修订、请求引用、状态与恢复；显式部署可用于预热或兼容工作流。
\item[Provider 管理] APPProvider/ProviderAdminPort 提供 stage、activate、drain、delete、
进度/检查点/输出报告；管理凭据与普通推理请求权限不是同一个授权入口。
\item[原生 host] NativeInferenceProvider 通过作用域注册句柄提供服务；
共享 ExecutionLease 状态与 target 路由绑定物理槽位，关闭/代次阻断晚到回调。
\item[扩展 SDK] registry、loader、bounded executor 和 observer 注册策略/适配器；
能力包含分割、放置、调度、准入、缓存、恢复、部署和执行调优等。
\end{description}

\subsection{可选后端与工具}
维护目录中有 ONNX、Qwen、YOLO adapter，Python 侧还保留 llama 等适配器与兼容入口；
模型导出、图检查、GUI、ops 和 experiment tooling 不属于通用 Core。
一个 API 导出不表示对应可选库、模型文件或 runtime 当前已安装。[D1, D8]

\subsection{持久记录}
部署修订、请求/attempt、计划摘要、Provider boot、grant、执行 lease 和结果引用需要一起定位，
不能只靠一个 deployment 名称恢复。恢复时应按 owner 的契约重新验证仍有效的运行资源。

\subsection{当前迁移状态}
原生 Provider host/注册/共享 lease 已有进度记录；Python facade 兼容迁移、
maintained caller 迁移及完整请求链仍未全部关闭。文档将 Python 控制路径和 C++ 组件同时列出，
避免把 Spec182 的未来单一实现方向覆盖当前仍实际存在的 owner。[D7]


\subsection{管理动作与请求动作的区别}
stage 准备某部署修订，activate 使其可供请求，drain 限制新增工作，delete 处理该部署资源；这些管理接口由 \code{APPProvider/ProviderAdminPort} 拥有，不能用 \code{NativeInferenceHandle.cancel} 代替。恢复定位需要部署修订、request/attempt、plan、Provider boot 和有效 lease/grant，而非只有可读部署名称。新增 adapter/strategy 必须被 registry/配置实际选择，步骤见 AC-23。

\newpage

\section{UAV：服务容器、接口与操作界面}
\lead{地面站与无人机都是可同时发起和提供服务的应用进程}
\code{UavDroneApp} 组合飞控适配、视频、遥测、任务与内嵌 Repo；
\code{UavGroundStationApp} 组合命令客户端、地图/任务界面、流订阅、记录播放和分析服务。
主体代码拆在 Drone/GroundStationServiceContainer 与 Window 的 include 文件中；
共享类型和命名位于 \code{shared/}。[U1, U2]

\begin{tabularx}{\linewidth}{@{}p{35mm}X@{}}
\toprule
服务组 & 当前名称/后缀及作用 \\
\midrule
飞行命令 & \code{/UAV/MAVLink/Execute}；受保护的 MAVLink 执行。\\
任务与状态 & \code{/UAV/Mission/Assign}、\code{/UAV/Telemetry/GetStatus}；任务分配与遥测。\\
摄像头 & \code{/UAV/Camera/GetFrame}、Provider 专属 \code{/UAV/Camera/Video}；取帧与开始/停止流。\\
录像与目录 & Provider 专属 \code{/UAV/Camera/Recording/Manifest}、\code{/UAV/Camera/Repo/Catalog}。\\
参数与诊断 & Provider 专属 Parameters、ParameterEdit、AnalyzeSnapshot 与 Preflight/Checklist。\\
识别与协作 & \code{/UAV/GS/ObjectDetection}、\code{/UAV/Incident/Analyze}。\\
操作权管理 & GS 的 OperatorAuthority/Lease、Revocation、Audit 服务。\\
\bottomrule
\end{tabularx}
名字可由 UavRuntimeConfig 配置；以上默认值不能代替部署配置和签名身份检查。

\subsection{界面与状态模型}
GUI 使用 typed TelemetryState、ReadinessState、FlightCommandState、MissionState、
MissionProgressState、VideoState、AnalyzeSnapshot 和 dashboard 数据，
区分无人机可达性、飞控 readiness、摄像头可用性、任务占用和视频会话。
参数浏览/修改、preflight、目录查询和诊断由服务调用支撑，不直接绕过 Core 访问远端设备。[U2, U3]

\callout{应用边界}{Core 传递 opaque MAVLink/应用 payload；飞行语义、地图与 GUI、媒体选择和检测模型均由 UAV 应用拥有。}


\subsection{界面如何使用状态}
例如命令按钮关联一次 invocation 的 requestId，飞控 readiness 来源于遥测，视频窗口关联一个当前 stream session；三者分别更新。关闭视频窗口不应被解释为 mission 已取消，取得服务响应也不证明飞控 COMMAND\_ACK 已到达。新增 UI 控件应订阅对应 typed state 并管理 Qt 回调生命周期，不直接重写设备状态。

\newpage

\section{UAV：飞控适配与操作权限}
\lead{网络授权、操作权和飞行安全属于三个责任域，实际门禁位置须分别核对}
\subsection{飞控数据路径}
GS 构造 MAVLink 命令，经 NDNSF 普通认证/token bootstrap 后对目标无人机使用 Targeted。
Drone 的后端负责向 mock、UDP 或 serial 飞控路径传递字节并读取 ACK/遥测；
mavlink-router 是部署连接方式之一，不是新的 Core 调用协议。[U2]

FlightControllerBackend 定义 sendMavlink、遥测、参数和 mission waypoint 接口。
UDP backend 包含连接、帧解析、heartbeat、命令 ACK 和 mission upload 交互，
并支持串口配置分支。GPS 信息来自飞控 MAVLink telemetry，而不是把任意本机 GPS 设备
直接当作 NDNSF 身份可信来源。

\subsection{三层权限与门禁}
\begin{enumerate}
\item Core 验证 User/Provider 权限、签名、token、ControllerVersion 与请求绑定。
\item UAV 的 OperatorAuthorityLease 关联操作员、目标、范围和有效时间；
issuer 提供申请、撤回、审计、持久化和续期逻辑。
\item GroundStation 的 flight/mission safety gate 使用新鲜遥测、readiness、
armed/landed 等应用状态判断动作是否允许；lease 不能替代飞控本身的保护。
\end{enumerate}
操作权 issuer 不等于 Core ServiceController；其 lease 修订也不是 Controller epoch。
UI 阻止命令、网络拒绝授权和飞控返回 command ACK 必须分别显示。[U2, U3]
当前 Drone Execute handler 解码 command/mavlink\_hex 后直接调用 backend，
其函数体没有操作 lease 或 readiness 的独立重验；不能把 GS 侧门禁当作 Drone 侧同样
强制执行的证据。AC-19 列出实际字段与 ACK 行为，TG-05 规划后端边界的明确校验。

\subsection{现有支持与限制}
mock/SITL、UDP、serial、遥测显示与参数操作为可运行入口；真实飞控、
fail-safe、长时间断连恢复和户外设备条件仍需对应运行验证。
代码具备硬件接口不表示本轮进行了飞行验证，也不等同于成熟飞控地面站的认证。


\subsection{沿失败位置排查}
命令尚未获得 Core 权限时停在服务入口；操作 lease 过期属于 UAV 操作权；设备链路断开或飞控拒绝属于 backend/飞控反馈。三类失败不可都显示为 timeout。AC-19 分开说明网络 Response、发送字节与 COMMAND\_ACK，同时明确当前 Fields 命令表未形成覆盖全部命令的静态 schema；类型化补齐属于 TG-05。

\newpage

\section{UAV：巡逻任务、补偿与事件协作}
\lead{长生命周期 mission 包含多个有限的 NDNSF invocation}
\subsection{巡逻分配}
地面站将 waypoint 集合确定性分区/聚类，形成任务 plan 与多个 part；
通过共享 \code{/UAV/Mission/Assign} 服务发现空闲 Drone。
忙碌 Drone 可以发负 ACK 并报告 mission slot 占用；应用根据候选选择，而不是用 payload
里的 Provider 名称绕过协商。[U2, U4]

\subsection{缺失部分的补偿}
patrol task 关联 taskId、attemptId、partId、assignedProvider、deadline、状态和 response digest。
某 part 超时或没有有效结果时，下一次调用只补偿缺失部分；每次仍有新的 Core requestId/token。
任务业务“完成”需要所有必要 part 的有效结果，而不是单次 invocation 有一个 Response。

\subsection{任务与 job 状态}
UavMissionSession 维护 planned、starting、active、degraded、compensating、recovering
到完成/取消/失败的任务状态，并限制历史大小。
每个 incident/job 另有 ACK 收集、ACK 关闭、计划提交、selected、evidence-ready、
executing、reporting 和 terminal 状态。[U3, U4]

\subsection{证据与终端归属}
证据使用生产者精确命名 Data，绑定 mission/incident、digest、时间与身份。
终端报告必须匹配当前 attempt、request、planDigest 和选中的 terminal owner；
同一 incident 已接收的终态摘要不会被另一份晚到报告替换。
恢复过程需要确认车辆与 stream 绑定，不能只载入一个旧 ledger 就继续发飞行命令。[U4]

\callout{与 Core 的分工}{Core 负责每次调用的认证与超时；UAV 负责 part 是否完成、哪些工作需要补偿、
跨调用的任务身份以及任务层最终完成。}


\subsection{两个分片的任务例子}
P0 已验证完成，P1 进入 missing 时，仅对 P1 开启补偿 attempt；P0 的结果继续保留。新 invocation 获得新 token/assignment，旧 P1 的迟到终端报告仍须匹配当前 attempt/plan/owner 才可接纳。恢复时先 reconcile 实际车辆与流，再允许新的业务动作。AC-20 解释 bool/reason、局部状态与真实飞行停止的差异。

\newpage

\section{UAV：视频、遥测流与录像}
\lead{小型控制请求建立会话，高频媒体通过命名 Data 传输}
\subsection{采集与呈现}
VideoPublisher 使用 GStreamer 或兼容管线，支持摄像头/文件/测试来源。
UavVideoPipeline 的 frame 结构携带 sessionEpoch、sourceFrameId、captureOriginNs、
codec PTS、codecConfigEpoch、keyFrame 与编码字节；有界最新帧队列控制积压。
readiness 报告真实来源及失败原因，fallback sample 不能当作真实摄像头采集证据。[U2, U5]

\subsection{网络流与自适应}
Provider 专属控制服务决定 start/stop；GS 订阅当前 live session，按 cursor/名称与元数据
预测、取回和重排 Data。旧会话迟到包不进入当前视频。
Core 提供通用窗口、名称映射、恢复与背压；UAV 根据帧类别、码率、画面时效和播放队列
做媒体策略，并记录采集、到达、解码与呈现的不同时间点。

\subsection{保密与留存}
视频会话拥有受保护的 descriptor、密钥/nonce 与 Data 绑定。原始已签名 packet 可以进入
Repo；留存清单记录确切 packet、缺口和会话身份，不通过重新签名抹去原发布者。
开始直播与是否开启 recording/retention 是不同控制条件；
录像完成后通过 manifest/catalog 发现，按名字获取、验证、解密并播放。[U3, U5]

\subsection{数据产品边界}
当前具体 Repo-backed 产品主要是相机录像/packet 留存。Mission image、telemetry log、
检测事件或报告虽可使用同样的 DataProductReference 模型，不能仅因类型存在就称所有产品均有完整采集和发布链。

\callout{时序边界}{网络收到、解码完成与实际显示不是同一事件；低延迟或稳定性结论需要真实采集/呈现记录，
不能用一个管线启动成功代替。}


\subsection{丢帧与录像完整性}
最新帧队列中有 10、11、12 时 popLatest 只交付 12，其余计入 dropped。push 返回 true 也可能因为丢弃旧帧才腾出空间；超大单帧直接返回 false。这个显示策略不能复制到完整 packet 录像，后者需 manifest 和 gap 记录。AC-21 同时给出遥测单位、admit 多字段结果及 capture→传输→解码→显示/录像分工。

\newpage

\section{UAV：目标检测与多视角识别}
\lead{有限识别 job 与长期任务分离，验证输入与结果的实际 owner}
\subsection{单视角辅助检测}
Drone 可调用 GS 的 ObjectDetection 服务，请求只携带 frame/drone/目标类别等紧凑元数据，
由 GS 对其已有的解码画面执行检测。YOLO worker 属于 UAV 应用侧计算工具；
此调用不能自动证明使用了完整 NDNSF-DI 分布式规划路径。[U2]

\subsection{多视角任务}
MultiViewRecognitionJob 持有不可变精确视图引用、模型 profile、目标、mission/job/attempt
和时间边界。UavDetectorProvider 验证获取的 Data，并通过
\code{IMultiViewRecognitionAlgorithm} 接口执行。
算法输出融合结果和 Provider-owned 标注引用；成功、视图不足、验证/关联失败、模型不支持、
计算失败、发布失败、超时与取消分别编码。[U6]

\subsection{真实模型与功能替身}
C++ 有 DeterministicMultiViewAlgorithm 用于显式功能契约检查。
真实 MVCNN-family ONNX 工具路径在 \code{tools/mvcnn_onnx_worker.py}：
检查注册模型摘要、视图摘要和输入形状，显式要求 CPUExecutionProvider；
images 与 viewMask 输入在图内联合处理，输出 logits 和 pooledFeatures。
常规多视角 job 要求不同 UAV 的足够视图；单视角配对 baseline 不能冒充多视角融合。[U6]

\subsection{不能扩大解释的结果}
C++ 算法注入接口存在不证明真实 ONNX worker 已在每个默认 GUI/Provider 路径中接入；
工具 fixture 的功能结果也不证明真实飞行准确率。
跨进程获取、候选竞争、单一终端 owner、结果/标注发布和失败终态需在对应网络夹具中分别验证。


\subsection{两视角输入到终端接受}
两架 Drone 提供同一事件窗口的精确视图引用，job 绑定这些输入与模型 profile。executeMultiViewJob 的算法输出随后仍要经过 coordinator 的 \code{job/attempt/terminalOwner} 检查；算法成功不等于业务结果已接受。默认确定性算法、真实 ONNX worker 和单帧 YOLO 是独立入口。扩展点和例子见 AC-22。

\newpage

\section{跨模块协作的完整路径}
\lead{用同一套 Core 机制连接领域任务与命名数据}
\subsection{UAV 录像路径}
GS 经 Core 认证调用 Drone 的视频控制；Drone 发布受保护媒体 Data，
可选内嵌 Repo 保存原始 packet 与留存 manifest。
GS 查询录像清单，按精确引用取回并验证/解密后播放。
UAV 定义媒体与保留语义；Repo 管字节/目录；Core 管调用、身份与传输。[U2, U5, R1]

\subsection{DI 模型意图路径}
应用提交模型和输入；Python 当前编排通过 Core 收集认证 ACK，
adapter/placement 决定图有效的角色分配；可信物化后经 Repo 发布制品，
再提交 Core Selection。Provider 原生运行器获取依赖、执行并发布输出，
terminal role 产生被调用者接收的结果。原生 requester 仍未接通同等全链路。[D2, D5, D7]

\subsection{UAV 分布式分析路径}
任务管理器记录 incident 和精确 evidence；通过 Core 协作收集分析 Provider，
固定角色/terminal owner，Provider 拉取并验证证据后返回报告及标注引用。
若某服务选择以 DI 实现计算，应通过明确的 DI 应用接口接入；
不能因同仓库有 DI 就把全部 UAV 检测自动画成 DI 调用。[U4, U6]

\begin{tabularx}{\linewidth}{@{}p{34mm}X@{}}
\toprule
共同约束 & 具体要求 \\
\midrule
控制/数据分离 & 服务请求传小型描述和引用；完整对象/媒体沿合适的数据面传输。\\
身份与摘要 & 调用者、发布者、Provider、对象、request、attempt、plan 等按各层契约绑定。\\
状态不混淆 & 存储已提交、模型已准备、资源已预留、飞行已执行、画面已显示是不同事实。\\
恢复有 owner & Core 清理调用；Repo 恢复已验证字节；DI 管执行/会话；UAV 管任务补偿。\\
\bottomrule
\end{tabularx}


\subsection{贯穿四模块的例子与实际接线边界}
一次 UAV 事件可以组织为：相机经 Core 命名流发布画面；录像路径将原 packet 与 manifest 存入 Repo；地面站以精确画面引用构造分析 job；分析服务验证输入并返回绑定 job/attempt 的结果。若选择接入 DI，DI 还要完成模型准备、ACK 规划、角色执行与终端验证。前半段现有 UAV/Repo 入口与可选 DI 集成必须分开核对，当前原生 requester 的 NOT\_READY 意味着不能把这条组合示例写成已经通过纯 C++ 端到端验收的调用链。

\newpage

\section{构建、配置与部署边界}
\lead{源码、宿主构建、容器候选和运行证据分别建立身份}
\begin{tabularx}{\linewidth}{@{}p{32mm}X@{}}
\toprule
部署形态 & 组件与前提 \\
\midrule
本地 Core 应用 & C++ runtime、ndn-cxx、NFD、ndn-svs、NAC-ABE、信任模式与私有身份环境。\\
Repo-node & Core 服务、SQLite/存储目录、热点缓存预算；大制品 backend 按所选物化方式配置。\\
UAV 应用 & Drone/GS 容器、GUI/媒体依赖、mock 或实际 MAVLink 连接、摄像头、可选 Repo。\\
DI Provider & Core、原生 DI 库、匹配模型 runner、ONNX Runtime/可选 CUDA、精确模型与制品引用。\\
MiniNDN & 多进程 NFD/路由/身份隔离、拓扑及实验脚本；不能用单进程调用替代网络证据。\\
Tiger/SIF & 封存源码、容器内原生构建、候选身份、Apptainer/Slurm、实际节点设备和执行记录。\\
\bottomrule
\end{tabularx}

\subsection{工具链一致性}
宿主原生构建按仓库规则使用 system compiler/binutils 和匹配的 Boost 头/库，
先完成 NAC-ABE 依赖闭包，再构建 Core 与 Python 扩展。
改变 C++ ABI 后，需要重新构建依赖对象并核对实际加载的库；
导入成功和 SONAME 不变都不能独立证明 ABI 一致。[O1]

\subsection{容器与集群}
Tiger 的唯一维护入口为 \path{Experiments/TigerCluster/}。
适配器、job/profile、images 和 results 按该目录契约组织；共享组件仍由 packaging 维护。
镜像原生组件须在容器 builder 内编译，不使用宿主 Python.h、venv 或宿主 .so 替代镜像依赖。
开发机负责代码/静态检查/测试，实验机负责 SIF 构建和 Tiger 实验。[O2]

\subsection{配置、身份和数据}
私钥、模型、数据库、缓存、运行日志与镜像不因文档生成进入版本库。
模型名、设备参数和服务路由需与实际被加载/注册的内容核对；
本轮没有启动 NFD/MiniNDN、构建原生依赖或向集群提交任务。


\subsection{可复现入口与本轮命令}
本轮文档命令从仓库根执行：python3 \code{Design/build-api-reference.py}；python3 \code{Design/build-behavior-coverage.py}；python3 \code{Design/refresh-snapshot.py}；python3 \code{Design/render-api-contracts.py} --target；python3 Design/build.py；python3 Design/verify.py 后接本次输出目录。原生编译参数按 \code{docs/native-build-parallelism.md} 与当前 Spec 的 quickstart/验证脚本，镜像按 \code{Experiments/TigerCluster/docs/sif-build.md}。文档更新不自行启动原生构建、SIF 或集群作业。

\newpage

\section{实现状态与验证口径}
\lead{把已有代码、接线状态和运行资格分开报告}
\begin{tabularx}{\linewidth}{@{}p{27mm}X@{}}
\toprule
范围 & 本版可支持的结论 \\
\midrule
Core & 已有通用调用、Targeted、协作、流、安全及在线授权组件；本文核对主要生产路径与状态分支。\\
Repo & 已有 C++ 对象存储/服务、Python 编排和制品 API；具体后端、目录操作及控制模式存在差异。\\
UAV & 已有飞控/GUI/任务/视频/Repo/识别组件和研究夹具；真实硬件和模型接线按入口分别确认。\\
DI & C++ PreparedModel/NativeInferenceClient 与 Provider 执行链已接通并由原生过程矩阵验证；旧 Python 编排、subinterpreter/wheel 与外部资格仍分别保留边界。\\
本轮文档 & 源码核对、PDF 构建与文本/版面检查；没有新协议测试或性能测量。\\
\bottomrule
\end{tabularx}

\subsection{测试入口}
Core 检查位于 \path{tests/unit-tests/} 与 \path{tests/integration-tests/}，
包含 controller-revocation、controller-version-refresh、invocation-stream、
stream-facade、opaque-selection 等主题。
Repo 有存储后端/制品传输 C++ 检查与 \code{test_spec164_*} Python 检查；
UAV 有 mission/collaboration/multiview 检查和 tools 运行夹具；
DI 有 \code{di-native-*} 及 Python/experiment 检查。[V1]

这些名字是继续验证的入口，不是本轮通过清单。测试文件存在、一个 fixture 运行成功、
完整网络结果和指定候选的硬件资格验收是不同证据等级。

\subsection{活跃实现}
当前 active Spec 为 \code{185-prepared-model-runtime}；其 C++ qualification、Python
binding 和 Design handoff 证据控制本轮状态。Spec184 的 Qwen3.6-27B、继承 negative/
retirement、I05、Python retirement 与 SIF/Tiger 资格仍由
\code{specs/184-native-di-closure/} 独立控制，本文不改写这些任务的验收状态。

\callout{失效快照的处理}{代码变化后，source-baseline 摘要可能与当前工作树不再一致。
这表示设计快照需要重新核对，而不是证明新代码有缺陷；目标文档也不能凭自动复制变成实现事实。}


\subsection{能力、入口与证据的对应}
Core 在线权限以 \code{ServiceController::grant/revoke} 和第 9–13 章分支为依据；Repo 精确目录以 RepoCore::catalogLookup 为依据；UAV 最新帧以 \code{BoundedLatestFrameQueue::push/popLatest} 为依据；DI 请求句柄以 NativeInferenceClient/PreparedModel 为依据。这些行为核对分别进入 AC-08、AC-13、AC-21、AC-14。生成/KV、原生过程、会话恢复和 Python facade 的当前证据分别链接 Spec185 B7/B8；Spec184 外部模型、retirement 与 SIF/Tiger 资格仍保持独立。Design/reviews 的修订记录只证明文档工作，不替代这些功能门。

\newpage

\section{源码索引：Core 与授权}
\lead{按入口定位，再用快照摘要确认正在阅读的版本}
下列路径以仓库根为起点；正文中的引用 ID 指向这里和下一节。
符号用于定位职责；详细文件清单与 SHA-256 见 \code{source-baseline.json}。

{\small
\begin{tabularx}{\linewidth}{@{}p{12mm}X@{}}
\toprule
编号 & 规范文件与关注点 \\
\midrule
C1 & \path{docs/architecture.md}；\path{docs/ndnsf-core-app-boundary.md}，总体边界；以源码修正旧叙述。\\
C2 & \path{ndn-service-framework/ServiceContainer.hpp}、\path{LocalServiceRegistry.hpp}；同目录 ServiceController/ServiceProvider。\\
C3 & Core 的 \path{ServiceUser.hpp}、\path{ServiceUser.cpp}、\path{ServiceProvider.cpp}、\path{NDNSFMessages.hpp}；调用与协作。\\
C4 & Core 的 \path{GenericSelectionTxnStore.hpp}、\path{GenericSelectionTxnStore.cpp}；不透明选择事务。\\
C5 & Core 的 ServiceUser/ServiceProvider、\path{RequestConfidentiality.cpp}、\path{HybridMessageCrypto.cpp}；大数据和加密。\\
C6 & Core 的 \path{Stream.hpp}、\path{Stream.cpp}、\path{StreamFacade.hpp}、\path{StreamFacade.cpp}；连续流、预测、FEC。\\
C7 & Core 的 \path{InvocationStream.hpp}、\path{InvocationStream.cpp}；调用事件流与终态。\\
C8 & Core 的 \path{ExecutionLease.hpp}、\path{NetworkTelemetry.hpp}、\path{TimelineTrace.hpp} 及 Provider；资源和观测。\\
\bottomrule
\end{tabularx}
}

\subsection{授权引用}
S1：ControllerVersion.cpp 的 isValid/compare。
S2：ServiceController.cpp 的 getPolicyStatus。
S3：同文件的 reissueAbePolicies、rotateAbeGenerationAndReissuePolicies、
reconcilePendingAbeRotation、advanceAuthorizationEpoch、revoke。
S4：同文件的 grant。
S5：User/Provider 的 applyPermissionResponse。
S6：installControllerStatus。
S7：invalidateControllerScopedCaches 与 Provider 的 refreshProviderPermissionsAfterAdvance。
S8：Controller 的 onPolicyStatusInterest、User/Provider 的 fetchPolicyStatusFromController/onPolicyStatusData。
S9：PolicyRefreshCoordinator.cpp。
S10：refreshNacDkeyForControllerStatus。
S11：scheduleControllerStatusRefresh 与周期配置。
S12：restorePersistedRuntimeStatuses。[均位于 Core 目录]

\newpage

\section{源码索引：Repo 与 DI}
\lead{文件清单包含规范入口；不以临时 compare/build 副本替代当前源文件}
{\small
\begin{tabularx}{\linewidth}{@{}p{12mm}X@{}}
\toprule
编号 & 路径（每组根目录见下文）及关注点 \\
\midrule
R1 & Repo \path{README_ch.md}、\path{include/ndnsf-distributed-repo/RepoCore.hpp}；对象与边界。\\
R2 & Repo \path{src/RepoNode.cpp}、\path{src/RepoClient.cpp}；服务与控制入口。\\
R3 & Repo \path{src/RepoTypes.cpp}；普通 SQLite/热点缓存后端。\\
R4 & Repo include 中 \path{RepoStoreBackend.hpp}、\path{FilesystemArtifactStore.hpp}、\path{ArtifactTransfer.hpp}、\path{ArtifactManifest.hpp}。\\
R5 & Repo \path{pythonWrapper/py_repoclient/artifact_api.py}；公开会话 API。\\
R6 & 同 Python 包 \path{network_artifact_backend.py}；真实网络发布、协作和 receipt 绑定。\\
R7 & 同 Python 包 \path{orchestration.py}、\path{persistence.py} 与 Repo \path{src/RepoCore.cpp}；目录、副本与后端。\\
D1 & DI Python \path{api/__init__.py}、\path{app_sdk/client.py}、\path{application.py}、\path{provider.py}。\\
D2 & DI Python \path{app_sdk/placement.py}、\path{sdk/placement.py}；ACK 驱动规划与 V3。\\
D3 & DI C++ \path{NativeRequestPreparation.cpp}、\path{NativeOfferAdmission.cpp}、\path{NativeV3Placement.cpp}。\\
D4 & DI C++ NativePlanSealer、NativeGrantClient/Verifier、ProtectedRuntime、NativeProtectedArtifactStore；封装和保护。\\
D5 & DI C++ NativeProviderSession/Runtime/Handler、ProviderRoleWorker、TensorBundleCodec、DependencyWaitScheduler、CollectiveRuntime。\\
D6 & DI C++ NativeEpochCoordinator、NativeConversationCoordinator、NativeStandaloneTokenizer；生成与会话。\\
D7 & DI C++ \path{PreparedModel.cpp}、\path{NativeInferenceClient.cpp}: prepared request/dispatchOperation；Spec185 tasks 与 evidence。\\
D8 & DI C++ NativeInferenceProvider；Python app\_sdk 与 sdk 的部署/扩展接口。\\
\bottomrule
\end{tabularx}
}
{\footnotesize
Repo 根为 \path{NDNSF-DistributedRepo/}；DI Python 根为
\path{NDNSF-DistributedInference/ndnsf_distributed_inference/}，DI C++ 根为
\path{NDNSF-DistributedInference/cpp/ndnsf-di/}。}


\subsection{生成与会话的准确文件入口}
D6 的生成自由函数位于 \code{NDNSF-DistributedInference/cpp/ndnsf-di/NativeEpochCoordinator.cpp}；\code{QwenGenerationSessionStateMachine} 的实现位于 \code{cpp/adapters/qwen/QwenGenerationSession.cpp}；会话日志位于 \code{cpp/ndnsf-di/NativeConversationCoordinator.cpp}；缓存状态机在 \code{NativeProviderRuntime.cpp} 的 ConversationStateStore。由 AC-18 的所属符号和签名定位，不把同名转发头当完整实现。

\newpage

\section{源码索引：UAV 与交付}
\lead{应用服务、硬件/模型工具与运行入口分别定位}
{\small
\begin{tabularx}{\linewidth}{@{}p{12mm}X@{}}
\toprule
编号 & 规范路径与关注点 \\
\midrule
U1 & UAV \path{README_ch.md}；\path{shared/UavNames.hpp}。\\
U2 & UAV \path{drone/DroneServiceContainer.inc.hpp}、\path{ground-station/GroundStationServiceContainer.inc.hpp} 及各 App/Window。\\
U3 & UAV \path{shared/UavProtocol.hpp}；状态与线协议类型。\\
U4 & UAV \path{shared/UavMissionSession.cpp}、\path{ground-station/UavIncidentCoordinator.cpp}、\path{drone/UavCollaborationParticipant.cpp}。\\
U5 & UAV \path{shared/UavVideoPipeline.hpp/.cpp} 及 Drone/GS 容器；媒体与留存。\\
U6 & UAV \path{shared/UavMultiViewRecognition.cpp}、\path{UavDetectorProvider.cpp}、\path{tools/mvcnn_onnx_worker.py}。\\
O1/O2 & \path{AGENTS.md}、\path{docs/native-build-parallelism.md}、\path{Experiments/TigerCluster/README.md} 与 \path{docs/sif-build.md}。\\
V1 & 根目录 \path{tests/}、\path{Experiments/}，各模块 examples/tools 及 active Spec evidence。\\
\bottomrule
\end{tabularx}
}

UAV 根为 \path{NDNSF-UAV-APP/}。O2 中的 \path{docs/sif-build.md} 相对
\path{Experiments/TigerCluster/}；其他交付路径相对仓库根。

\subsection{模块维护文件清单}
\code{module-inventory.json} 按 Core、UAV、DI、Repo 以及公共绑定/支持目录记录受版本库维护的
文件路径、类型、大小及源码摘要；模型、原始日志和构建目录不因文档任务被复制进 Design。
重点核对的规范文件另列于 \code{source-baseline.json}。

\subsection{源码检索口径}
CodeGraph 用于定位规范符号和调用者；出现 compare-old/compare-new、staging 或 build
副本时，以仓库维护的原始路径为准。章节中的运行限制来自具体实现分支，
不是根据类名或 README 的目标描述推测。


\subsection{索引的使用范围}
UAV 高频回调和硬件行为常在 .inc.hpp 与 .cpp 中，公开头文件清单不能代替这些实现。video queue 查 \code{shared/UavVideoPipeline.cpp}，mission 查 \code{shared/UavMissionSession.cpp}，实际命令解析查 \code{drone/DroneServiceContainer.inc.hpp}。module-inventory.json 是 R0 历史盘点，当前新增/删除检测以 \code{design_state.source_files} 和 source-baseline 为准。

\newpage

\section{文档维护与后续目标变更}
\lead{当前设计随已核对代码更新；目标设计随用户明确决定更新}
\begin{enumerate}
\item 当前正文保存在 \code{current-content.tex}，目标正文保存在 \code{target-content.tex}。
首版完全相同；以后两份独立维护，不用共享正文把目标修改自动传播到当前设计。
\item 修改目标时，记录变更动机、涉及模块、与当前行为的差异、依赖和验证条件。
未决定的方案保持待讨论，不把想法当作已接受契约。
\item 实现落地后，重新核对生产路径、状态/安全分支和实际证据，
更新当前正文与源码快照；不能仅因为目标 PDF 已更新就宣称功能完成。
\item 运行 \code{python3 Design/build.py} 重建两份 PDF，检查文本、分页、字体、表格和图。
当前与目标按各自源码身份检查；TG-01 至 TG-05 的有意差异必须保留。
\item 更新 \code{CHANGELOG.md}、\code{validation.md} 和相关 Spec 进度。
本目录完整纳入 Git；每个 Spec 的设计影响记入 \code{spec-design-changes.md}，与正文和 PDF 同步提交。
\end{enumerate}

\subsection{覆盖清单与源码快照}
\code{coverage-matrix.md} 记录四模块各功能域对应章节与主要入口；
\code{module-inventory.json} 保留 R0 历史盘点；新增/删除由当前源码范围检查发现。
文件被列入清单不表示每一行都完成审计；正文区分主要流程源码核对与配置/工具入口登记。

\code{source-baseline.json} 记录重点核对文件的摘要、采样时提交和工作树差异。
本地源码快照与 PDF 构建日志位于唯一的 \code{.codex-tmp/} 目录，
供核对文档所描述的时点；它们不授予任何运行资格。

\subsection{本文未引入的行为变化}
本版不新增 SVS epoch 公告、不改变 grant 的本地失效范围、
不补接原生 DI 请求链，也不扩大 Repo/UAV 现有接口能力。
目标 TG-01 至 TG-05 已登记改进方向；下一步沿 Spec182 推进 TG-02，并逐项定稿其余目标的接口和验收条件。

\callout{交付标准}{两份完整中文 PDF、可编辑源文件、覆盖/源码记录和验证记录一起维护。
文档完成与代码实现完成是两个独立状态。}

\subsection{章节的人工接受条件}
读者应能从每个关键入口说明参数从哪里来、下一步调用什么、怎样判定完成以及失败保留在哪一层。签名生成和 PDF 校验无法自动证明这些问题已回答。修改需连同 API、正文、目标差异、Spec 记录、PDF 和验证证据形成一个本地提交；未定义的字段/线程保证明确登记，禁止用增加页数替代行为契约。
